% ADR-15 — se incluye desde DescripcionArquitecturaSoftware.tex con \input,
% igual que ADR-14, al que complementa.

\subsection*{ADR-15 --- Correcciones a la vista de contenedores y decisiones que el diagrama dejaba abiertas}

\textbf{Estado:} \textit{aceptado; complementa a ADR-14}, no lo supera. El estilo
híbrido, los catorce servicios y el criterio de extracción no cambian. Lo que
cambia son \textbf{cuatro relaciones} del diagrama de contenedores y
\textbf{seis puntos} que ADR-14 no fijaba.

\textbf{Contexto:} la revisión del diagrama de contenedores de ADR-14 lo
contrastó contra una especificación escrita a partir de él y encontró dos tipos
de problema.

\begin{itemize}
  \item \textbf{Relaciones que contradicen el propio diseño.} Acceso QR llamaba
        de forma síncrona a Ventas para gastar cada QR, aunque ADR-14 lo separa
        justamente porque su pico ---la puerta--- ocurre en otro momento que el
        de la venta: con esa llamada, una Ventas lenta atascaba la puerta. El
        \textit{webhook} de la pasarela llegaba directo a Pagos, saltándose el
        punto único de entrada. Y el gateway le preguntaba a Autenticación por
        cada petición, lo que ponía a quien emite los tokens en el camino de
        todo el tráfico.
  \item \textbf{Decisiones sin tomar.} El diagrama no dice con qué se construyen
        los servicios, qué broker ni qué motor de búsqueda se usan, cuántas
        réplicas mínimas llevan los microservicios que «escalan solos», cómo se
        invalida el QR anterior en una reventa, cómo se avisa de un evento
        cancelado, ni cómo se garantiza que un cambio guardado llegue al broker.
\end{itemize}

\textbf{Decisión --- relaciones corregidas:}

\begin{enumerate}
  \item \textbf{El API Gateway valida el JWT localmente, con la clave pública},
        y autoriza por el rol del token. No llama a Autenticación en cada
        petición: Autenticación solo emite tokens (\textit{login}, renovación).
        Los servicios siguen verificando el token por su cuenta y comprobando en
        Redis que la sesión no esté cerrada (RNF-06, ADR-11).
  \item \textbf{El gateway exige el turno de la Cola virtual en la ruta de
        compra} y lo gasta antes de reenviar la petición a Ventas. Ventas no
        llama a la Cola: solo recibe a cuantos puede atender. El turno es de un
        solo uso y solo vale para la cuenta que lo pidió.
  \item \textbf{Acceso QR no llama a Ventas.} Mantiene su propia copia del estado
        de cada QR, alimentada por \texttt{CompraConfirmada} (activa los QR) y
        por un evento nuevo, \texttt{EntradaActualizada}, que Ventas publica
        cuando una entrada sale a reventa, vuelve, cambia de dueño o se anula.
        La validación pasa a ser \textbf{una transacción local} ---QR, aforo y
        registro del ingreso--- en lugar de una Saga con compensación remota.
        En sentido contrario, Acceso QR publica \texttt{IngresoRegistrado} y
        Ventas marca la entrada como usada: sin eso permitiría reembolsar
        (CU-003D) o revender una entrada con la que ya se entró.
  \item \textbf{El \textit{webhook} de la pasarela entra por el gateway}, por la
        única ruta pública sin JWT; Pagos lo autentica con la firma HMAC del
        cuerpo. Con él, Pagos publica \texttt{PagoAprobado} /
        \texttt{PagoRechazado}, y un cobro que quedó indeterminado por un
        \textit{timeout} se resuelve cuando la pasarela confirma.
\end{enumerate}

\textbf{Decisión --- lo que el diagrama dejaba abierto:}

\begin{enumerate}
  \setcounter{enumi}{4}
  \item \textbf{Tecnologías.} Backend en \textbf{NestJS} (ya decidido en ADR-09);
        bases en \textbf{PostgreSQL} (ADR-06); broker \textbf{RabbitMQ} (ADR-10)
        con colas \textit{quorum}, que son su forma de «tópicos replicados»; e
        índice de búsqueda sobre una \textbf{instancia PostgreSQL dedicada con
        búsqueda de texto completo} (\texttt{tsvector} con índice GIN y
        \texttt{pg\_trgm} para errores de escritura). Se prefiere a OpenSearch
        porque los filtros de CU-005 ---categoría, fecha, ciudad, artista--- no
        necesitan relevancia ni facetas, y así el índice no añade una
        tecnología más que operar. Se revisa si la búsqueda empieza a necesitar
        ordenar por relevancia.
  \item \textbf{Réplicas mínimas:} las de la vista física
        (§\ref{sec:fisica}), que el diagrama no repetía. Ventas \textbf{3},
        igual que Cola virtual y Catálogo, porque los tres forman el camino de
        compra; Pagos y Autenticación \textbf{3}, porque si caen degradan muchos
        casos de uso a la vez; Acceso QR \textbf{2}, porque su carga se
        concentra en las horas de ingreso y escala por validaciones por segundo.
  \item \textbf{Transactional Outbox} para todo evento que un servicio publique
        como consecuencia de escribir en su base: el evento se guarda en una
        tabla de salida \textit{en la misma transacción} que el cambio, y un
        publicador lo saca al broker después. Se elige frente a «publicar
        después del commit» porque ese hueco pierde el evento si el broker o el
        proceso caen, y sin \texttt{CompraConfirmada} una entrada pagada no
        entraría por la puerta. La entrega es «al menos una vez»: cada evento
        lleva un \texttt{id} y una \texttt{secuencia} para descartar repetidos y
        desordenados.
  \item \textbf{\texttt{EventoCancelado}}, publicado por Eventos y recintos y
        consumido por Catálogo (lo retira), Ventas (cierra la venta), Acceso QR
        (rechaza en la puerta) y Notificaciones (avisa a los compradores). Con
        solo \texttt{EventoPublicado}, cancelar no tenía cómo llegar a nadie.
\end{enumerate}

Con esto el diagrama tiene \textbf{27 contenedores}: 4 \textit{frontends}, el
gateway, 7 microservicios, 7 servicios de dominio, 5 almacenes, el broker y 2
sistemas externos. La especificación de la revisión decía 34, pero su propia
tabla suma 27.

\textbf{Alternativas consideradas:}
\begin{itemize}
  \item \textbf{Mantener la llamada de Acceso QR a Ventas} y protegerla con un
        \textit{circuit breaker}. Acota el daño pero no lo quita: con el
        circuito abierto la puerta deja de validar, que es justo lo que ADR-14
        quería evitar.
  \item \textbf{Que Ventas gaste el turno llamando a la Cola.} Funciona, pero
        añade una dependencia síncrona más en el servicio que más carga
        soporta, y una flecha que el diagrama no tiene.
  \item \textbf{Validar el JWT en cada servicio y no en el gateway.} Ya se hace
        (RNF-06); hacerlo también en el gateway rechaza lo inválido antes de
        consumir un salto de red y una réplica.
\end{itemize}

\textbf{Atributo de calidad:} \textbf{Disponibilidad} ---la puerta ya no
depende de Ventas, ni el tráfico de Autenticación---; \textbf{Rendimiento},
porque una validación de QR y una autenticación dejan de costar un salto de
red; \textbf{Consistencia}, por el Outbox y por la validación atómica en la
puerta; y \textbf{Seguridad}, porque el turno solo vale para su dueño y el
\textit{webhook} solo se acepta firmado.

\textbf{Consecuencias:}
\begin{itemize}
  \item \textbf{Consistencia eventual en la puerta.} Una entrada recién comprada
        tarda en poder entrar lo que tarda su evento en llegar ---medio segundo
        en condiciones normales---. Para una puerta que abre horas después de la
        compra es irrelevante.
  \item \textbf{RNF-02 se cumple por orden, no por atomicidad.} Antes, el QR
        viejo y el nuevo se cambiaban en el mismo \texttt{UPDATE} que consultaba
        la puerta. Ahora la puerta deja de aceptar la entrada cuando sale a
        reventa ---mucho antes de transferirla---, así que el QR viejo ya no
        vale cuando llega el nuevo. Queda una ventana: entre que Ventas marca la
        entrada \texttt{EN\_REVENTA} y la puerta se entera, el vendedor podría
        entrar con ella. Si después se transfiere, Acceso QR conserva el estado
        \texttt{USADA} y registra una alerta para resolverla con el comprador.
  \item \textbf{Ventas ya no sabe quién entró.} \texttt{USADA} vive solo en
        Acceso QR. No afecta a la cancelación ni a la reventa, que cierran antes
        de que empiece el evento.
  \item \textbf{Una sesión cerrada pasa el gateway} y la rechaza el servicio de
        destino. La garantía es la misma; cambia quién responde el 401.
  \item \textbf{Un turno se gasta aunque la reserva falle} ---por ejemplo, si la
        localidad se agotó---. Es lo correcto para el aforo de la cola, pero
        obliga al comprador a volver a la fila si quiere intentar otra
        localidad.
\end{itemize}

\textbf{Estado de implementación.} Ya en el código: la validación local del JWT,
la autorización por rol y la exigencia del turno en el gateway
(\texttt{App/gateway/autenticacion.js}); la ruta pública del \textit{webhook}
con firma HMAC y la publicación de \texttt{PagoAprobado} /
\texttt{PagoRechazado} en Pagos; el Outbox en Ventas para
\texttt{CompraConfirmada} y \texttt{EntradaActualizada} y en Acceso QR para
\texttt{IngresoRegistrado}; y la copia local de los QR en Acceso QR, que ya no
llama a Ventas. Pendiente, porque depende de
servicios que todavía no existen: \texttt{EventoPublicado} y
\texttt{EventoCancelado} (Eventos y recintos), el consumo de
\texttt{PagoAprobado} en Ventas para cerrar la Saga (fase E de
\texttt{MIGRACION-ADR-14.md}), el Outbox en los demás publicadores, y el índice
de búsqueda dedicado del Catálogo.
